\chapter{Conclusiones}\label{chap:conclusiones}

Este capítulo presenta las conclusiones del proyecto en relación con sus cuatro objetivos específicos. Distingue los objetivos completados de aquellos que permanecen en desarrollo y limita las afirmaciones a los resultados documentados en el Capítulo~\ref{chap:resultados}.

\section*{Objetivo 1: caracterización de conjuntos de datos}

El objetivo de caracterizar conjuntos de imágenes dermatológicas pertinentes para las lesiones cutáneas de interés en Santander presenta avances, pero aún no se considera cerrado. Se describieron HAM10000, ISIC2019, ISIC2020 y PAD-UFES-20, incluidas sus clases diagnósticas, la disponibilidad de fototipo clínico y las limitaciones de los valores de tono calculados a partir de las imágenes (Sección~\ref{sec:conjuntos-datos-resultados} y Sección~\ref{sec:caracterizacion-tono}).

La revisión de literatura regional, reportes del Instituto del Cáncer del HIC y boletines del SIVIGILA permitió documentar quince patologías cutáneas de relevancia clínica en Santander y organizarlas en cinco niveles de prioridad (Tabla~\ref{tab:patologias-santander}). La caracterización también mostró que las patologías inflamatorias, parasitarias e infecciosas tropicales identificadas en la región no están cubiertas por los cuatro conjuntos descritos.

Los fototipos III y IV se consideran de interés para el proyecto, pero las fuentes locales revisadas no permiten afirmar que sean predominantes en Santander ni asignarles una proporción poblacional. Para cerrar este objetivo se requiere completar la trazabilidad de los criterios de selección y exclusión, las licencias, los grupos de lesión, las particiones y las fuentes de las afirmaciones epidemiológicas.

\section*{Objetivo 2: diseño del algoritmo de clasificación autosupervisado}

El objetivo de diseñar un algoritmo de clasificación de lesiones cutáneas basado en aprendizaje autosupervisado se considera logrado. El proyecto documentó tres líneas técnicas: el clasificador SimSiam anclado sobre un \emph{backbone} DINOv2 ViT-B/14 con adaptación de dominio; el módulo complementario de segmentación sin etiquetas, basado en MoCo v2, CCAM, pseudo-máscaras y un U-Net estudiante con refinamiento iterativo; y el procedimiento determinístico de transformación y estimación del tono de piel en el espacio CIELAB. Los algoritmos se describen en el Capítulo~\ref{chap:metodologia}.

La documentación reúne las configuraciones metodológicas registradas para estas líneas. En D2 y D3, sin embargo, la ausencia de un valor válido en \texttt{data\_revision} limita la trazabilidad exacta de las particiones utilizadas. Esta limitación debe considerarse al interpretar y reproducir esos resultados.

\section*{Objetivo 3: evaluación del desempeño con métricas de la literatura}

El objetivo de evaluar el desempeño mediante métricas utilizadas en la literatura se considera logrado dentro de las evaluaciones documentadas. Los resultados muestran diferencias según la tarea, el conjunto de datos y el protocolo aplicado.

La integración de segmentación en la clasificación tampoco demostró una mejora general. En la evaluación propia de segmentación, S4-A2 obtuvo un Dice de 0,7221 en ISIC2018 y 0,7949 en HAM10000, y superó a \texttt{student\_r3} entre las dos variantes evaluadas. Este resultado caracteriza su desempeño como segmentador, pero no demuestra que el uso de sus máscaras mejore la clasificación.

La evaluación del estimador ITA produjo exactitudes inferiores a la línea base por mayoría. En Fitzpatrick17k-C, el ITA obtuvo entre 28,5\,\% y 30,8\,\% de exactitud, frente a 32,9\,\% de la línea base. En PAD-UFES-20, la exactitud fue 21,74\,\% con imagen completa (n = 1.265; línea base de 59,37\,\%) y 22,88\,\% con el fondo enmascarado por Otsu (n = 1.189; línea base de 59,80\,\%). Ambas comparaciones se calcularon sobre los mismos casos evaluables para cada condición. La cifra publicada de 22,9\,\% frente a 51,3\,\% corresponde a otro protocolo, no a la evaluación del proyecto \parencite{alipour2026italimitations}. En consecuencia, el ITA no debe interpretarse como una estimación validada del fototipo clínico. Esta evaluación tampoco demuestra el efecto de la transformación de tono: aún se requieren mediciones antes y después sobre las mismas imágenes, con idéntico preprocesamiento.

En la clasificación congelada de HAM10000, B71 obtuvo intervalos pareados que excluyeron cero frente a ocho líneas base. No se aplicó una corrección global por multiplicidad, por lo que este patrón se interpreta dentro de la comparación evaluada. El ajuste fino actualiza todos los pesos con etiquetas y utiliza particiones distintas de la comparación congelada. Sus resultados no permiten atribuir diferencias únicamente al régimen de entrenamiento ni establecer una superioridad mediante una comparación pareada con B71. Por tanto, los resultados no establecen una superioridad general de B71; muestran el desempeño de una representación aprendida mediante SSL bajo el protocolo congelado evaluado.

En ISIC2020, los resultados de F1 macro y sensibilidad de melanoma apuntan en direcciones distintas. B71 detectó 39 de 122 melanomas (\textit{recall} 0,3197), mientras que AlexNet ajustado detectó 86 de 122 (\textit{recall} 0,7049). Aunque B71 obtuvo mayor F1 macro frente a siete de ocho líneas base ajustadas en el análisis post hoc, ese resultado no implica una mayor sensibilidad para melanoma. El proyecto no establece un criterio clínico para priorizar una de estas métricas en la selección final.

La adaptación de dominio con PAD-UFES-20 mostró una mejora piloto de D3 frente al control de cómputo D2 en F1 ($\Delta=+0{,}0816$). D3, no obstante, permaneció por debajo del \textit{encoder} público D0. La evaluación se basó en una semilla y en pocos casos de melanoma, por lo que no permite concluir que la adaptación al dominio clínico esté resuelta.

\section*{Objetivo 4: prototipo de aplicación de apoyo al análisis de imágenes}

El prototipo de apoyo al análisis de imágenes dermatológicas, previsto en el objetivo 4, aún no se ha implementado ni desplegado. Su desarrollo continúa pendiente mientras el equipo experimenta para seleccionar el modelo y la configuración que se utilizarían en esa etapa.

\section*{Limitaciones transversales}

Los resultados deben interpretarse considerando cuatro límites. Primero, los registros de D2 y D3 no incluyen un valor válido de \texttt{data\_revision}, lo que reduce la trazabilidad de sus particiones. Segundo, las fuentes revisadas no explican por qué se definieron exactamente tres rondas de autoentrenamiento para el segmentador. Tercero, los pocos casos de fototipos V y VI en PAD-UFES-20 no permiten conclusiones sobre esos grupos. Cuarto, ninguna evaluación de este proyecto constituye validación clínica ni reemplaza la valoración profesional.

\section{Recomendaciones y trabajo futuro}

\subsection{Mejoras metodológicas}

Completar los metadatos de procedencia de D2 y D3, incluidos \texttt{data\_revision}, las particiones, los hashes y la versión del \textit{pool}, para facilitar la revisión y reproducción de los resultados. En la ablación F4, evaluar por separado el recorte de la lesión y los descriptores de forma y color para identificar la contribución de cada componente. También se recomienda comparar configuraciones con distinto número de rondas de autoentrenamiento y documentar el criterio usado para seleccionar una variante. Para completar el objetivo 1, registrar los criterios de selección y exclusión, las licencias, los grupos de lesión, los hashes de partición y la fuente primaria de cada afirmación epidemiológica.

\subsection{Validación externa}

Evaluar el ITA frente a anotaciones expertas en un conjunto de mayor tamaño, dado que su desempeño actual quedó por debajo de la línea base por mayoría en Fitzpatrick17k-C y PAD-UFES-20. Ampliar la evaluación de adaptación en PAD-UFES-20, donde la partición \texttt{target\_test} incluye pocos melanomas y el resultado disponible corresponde a una semilla. Definir con especialistas qué criterio clínico debe orientar la selección del modelo, considerando la diferencia observada entre F1 macro y sensibilidad para melanoma en ISIC2020.

\subsection{Trabajo futuro clínico}

Desarrollar el prototipo académico de apoyo a la revisión de imágenes contemplado en el objetivo 4 cuando se hayan seleccionado el modelo y la configuración para esa etapa. Antes de considerar un uso de apoyo diagnóstico, evaluar el sistema con dermatólogos y datos clínicos de Santander. También se puede estudiar la cobertura de patologías regionales que no están representadas en los conjuntos públicos, entre ellas las inflamatorias, las infecciosas tropicales y la lepra, documentadas en la Sección~\ref{sec:perfil-epidemiologico-santander}.
